\documentclass[ ../main.tex]{subfiles}
\providecommand{\mainx}{..}
\begin{document}

\section{Composition of approximate maps}
\label{sec:composition}

If approximate maps are the morphisms of a category of approximate
computations, then composition must be well-defined and its error
behaviour must be predictable~\cite{cover2006}.  This section establishes that the
composition of two independent approximate maps is itself an approximate
map, derives closed-form expressions for the false positive and false
negative rates of the composition, and examines the passage from
first-order to higher-order models that composition induces.

\subsection{Composition setup}
\label{sec:comp-setup}

Let $\Fun{f} \colon \Set{X} \pfun \Set{Y}$ be a partial function with
domain of definition $\Dod(\Fun{f}) \subseteq \Set{X}$, and let
$\Fun{g} \colon \Set{Y} \pfun \Set{Z}$ be a partial function with
domain of definition $\Dod(\Fun{g}) \subseteq \Set{Y}$.
Lifting to the $\Maybe{}$ type, we write
$\Fun{f} \colon \Set{X} \to \Maybe{\Set{Y}}$ and
$\Fun{g} \colon \Set{Y} \to \Maybe{\Set{Z}}$.

\begin{assumption}[Composability]
\label{asm:composability}
For every $x \in \Dod(\Fun{f})$,
$\Fun{f}(x) \in \Dod(\Fun{g})$.
\end{assumption}

Under \cref{asm:composability} the ordinary composition
$\Fun{g} \circ \Fun{f} \colon \Set{X} \pfun \Set{Z}$ is well-defined
with $\Dod(\Fun{g} \circ \Fun{f}) = \Dod(\Fun{f})$, and for every
$x \in \Dod(\Fun{f})$,
$(\Fun{g} \circ \Fun{f})(x) = \Fun{g}(\Fun{f}(x))$.

Now let
$\APFun{f}[\fprate_1][\fnrate_1] \colon \Set{X} \to \Maybe{\Set{Y}}$
and
$\APFun{g}[\fprate_2][\fnrate_2] \colon \Set{Y} \to \Maybe{\Set{Z}}$
be independent approximate maps of $\Fun{f}$ and $\Fun{g}$
respectively, each using the Maybe-type model: for every
$x \in \Dod(\Fun{f})$, $\APFun{f}[\fprate_1][\fnrate_1](x)$ returns
either the correct value $\Fun{f}(x)$ or $\Nothing$ (and likewise for
$\APFun{g}[\fprate_2][\fnrate_2]$).
Independence means that the random outcomes of
$\APFun{f}[\fprate_1][\fnrate_1]$ and
$\APFun{g}[\fprate_2][\fnrate_2]$ are determined by independent
sources of randomness.

\begin{definition}[Composition of approximate maps]
\label{def:approx-comp}
The \emph{composition}
$\APFun{g}[\fprate_2][\fnrate_2] \circ
 \APFun{f}[\fprate_1][\fnrate_1]
 \colon \Set{X} \to \Maybe{\Set{Z}}$
is the random function defined, for each $x \in \Set{X}$, by
\begin{equation}
\label{eq:comp-def}
\left(\APFun{g}[\fprate_2][\fnrate_2] \circ
      \APFun{f}[\fprate_1][\fnrate_1]\right)(x)
\;\coloneqq\;
\begin{cases}
\APFun{g}[\fprate_2][\fnrate_2](y)
  & \text{if } \APFun{f}[\fprate_1][\fnrate_1](x) = y \in \Set{Y},\\[4pt]
\Nothing
  & \text{if } \APFun{f}[\fprate_1][\fnrate_1](x) = \Nothing.
\end{cases}
\end{equation}
\end{definition}

\begin{lemma}[Closure under composition]
\label{prop:closure}
Under \cref{asm:composability}, the composition
$\APFun{g}[\fprate_2][\fnrate_2] \circ
 \APFun{f}[\fprate_1][\fnrate_1]$
is an approximate map of $\Fun{g} \circ \Fun{f}$.
\end{lemma}
\begin{proof}
Let $x \in \Dod(\Fun{g} \circ \Fun{f}) = \Dod(\Fun{f})$.
With probability $1 - \fnrate_1$,
$\APFun{f}[\fprate_1][\fnrate_1](x) = \Fun{f}(x)$.
By \cref{asm:composability},
$\Fun{f}(x) \in \Dod(\Fun{g})$, so with probability
$1 - \fnrate_2$,
$\APFun{g}[\fprate_2][\fnrate_2](\Fun{f}(x)) = \Fun{g}(\Fun{f}(x))
= (\Fun{g} \circ \Fun{f})(x)$.
By independence, the probability that the composition returns the
correct value is $(1 - \fnrate_1)(1 - \fnrate_2) > 0$.

Conversely, let $x \notin \Dod(\Fun{f})$.
With probability $1 - \fprate_1$,
$\APFun{f}[\fprate_1][\fnrate_1](x) = \Nothing$, so the composition
returns \Nothing.
Hence $\Prob{(\APFun{g}[\fprate_2][\fnrate_2] \circ
\APFun{f}[\fprate_1][\fnrate_1])(x) = \Nothing}
\geq 1 - \fprate_1 > 0$.
\end{proof}

% -------------------------------------------------------------------
\subsection{Error propagation}
\label{sec:error-prop}

We now derive the false negative and false positive rates of the
composed map.  Throughout, we write
$\APFun{h} \coloneqq
\APFun{g}[\fprate_2][\fnrate_2] \circ
\APFun{f}[\fprate_1][\fnrate_1]$
for the composition, and
$\Fun{h} \coloneqq \Fun{g} \circ \Fun{f}$ for the objective
composition.  The domain of definition is
$\Dod(\Fun{h}) = \Dod(\Fun{f})$.

\subsubsection{False negative rate}
A false negative in the composition occurs when
$x \in \Dod(\Fun{h})$ but $\APFun{h}(x) = \Nothing$.

\begin{theorem}[Composition false negative rate]
\label{thm:comp-fnr}
Let $\APFun{f}[\fprate_1][\fnrate_1]$ and
$\APFun{g}[\fprate_2][\fnrate_2]$ be independent approximate maps
under the Maybe-type model (i.e., for positive inputs, each map
returns either the correct value or $\Nothing$), satisfying
\cref{asm:composability}.
Then for every $x \in \Dod(\Fun{h})$,
\begin{equation}
\label{eq:comp-fnr}
\fnrate_{\Fun{g} \circ \Fun{f}}
\;=\; 1 - (1 - \fnrate_1)(1 - \fnrate_2)
\;=\; \fnrate_1 + (1 - \fnrate_1)\,\fnrate_2\,.
\end{equation}
\end{theorem}

\begin{proof}
Fix $x \in \Dod(\Fun{h}) = \Dod(\Fun{f})$.  The composition returns
\Nothing if and only if one of the following mutually exclusive events
occurs:
\begin{enumerate}
    \item[(i)] $\APFun{f}[\fprate_1][\fnrate_1](x) = \Nothing$ (false
    negative from $\APFun{f}$), which has probability $\fnrate_1$; or
    \item[(ii)] $\APFun{f}[\fprate_1][\fnrate_1](x) = \Fun{f}(x)$
    (correct output from $\APFun{f}$) but
    $\APFun{g}[\fprate_2][\fnrate_2](\Fun{f}(x)) = \Nothing$ (false
    negative from $\APFun{g}$).
\end{enumerate}
Since $\APFun{f}[\fprate_1][\fnrate_1](x)$ is either $\Fun{f}(x)$ or
$\Nothing$ for $x \in \Dod(\Fun{f})$ (by the Maybe-type model),
these cases are exhaustive.
By \cref{asm:composability}, $\Fun{f}(x) \in \Dod(\Fun{g})$, so
in case~(ii) the probability of a false negative from $\APFun{g}$ is
$\fnrate_2$.  By independence:
\begin{align}
\fnrate_{\Fun{g} \circ \Fun{f}}
  &= \Prob{\APFun{h}(x) = \Nothing \Given x \in \Dod(\Fun{h})}
    \notag \\
  &= \Prob{\APFun{f}[\fprate_1][\fnrate_1](x) = \Nothing}
   + \Prob{\APFun{f}[\fprate_1][\fnrate_1](x) = \Fun{f}(x)} \cdot
     \Prob{\APFun{g}[\fprate_2][\fnrate_2](\Fun{f}(x)) = \Nothing}
    \notag \\
  &= \fnrate_1 + (1-\fnrate_1)\,\fnrate_2\,.
    \label{eq:fnr-derivation}
\end{align}
Rearranging, $\fnrate_{\Fun{g} \circ \Fun{f}} =
1 - (1-\fnrate_1)(1-\fnrate_2)$.
\end{proof}

\begin{remark}
The true positive rate of the
composition is the product of the individual true positive rates:
$\tprate_{\Fun{g} \circ \Fun{f}} = \tprate_1 \, \tprate_2$.
\end{remark}

\subsubsection{False positive rate}
The false positive analysis is more delicate because the false positive
value produced by $\APFun{f}[\fprate_1][\fnrate_1]$ may or may not
land in $\Dod(\Fun{g})$, and $\APFun{g}[\fprate_2][\fnrate_2]$ treats
these two sub-cases differently.

\begin{definition}[False-positive domain overlap]
\label{def:rho}
Let $\rho$ denote the probability that a false positive value
produced by $\APFun{f}[\fprate_1][\fnrate_1]$ falls in
$\Dod(\Fun{g})$:
\begin{equation}
\label{eq:rho-def}
\rho
\;\coloneqq\;
\Prob{y \in \Dod(\Fun{g}) \Given
      y \text{ is a false positive output of }
      \APFun{f}[\fprate_1][\fnrate_1]}\,.
\end{equation}
\end{definition}

The parameter $\rho$ captures the
interaction between the false positive distribution of $\APFun{f}$ and
the domain structure of $\Fun{g}$.
If $\APFun{f}$ produces false positives uniformly on $\Set{Y}$, then
$\rho = \Card{\Dod(\Fun{g})}/\Card{\Set{Y}}$, but in general $\rho$
depends on the implementation.

\begin{theorem}[Composition false positive rate]
\label{thm:comp-fpr}
Under the hypotheses of \cref{thm:comp-fnr}, for every
$x \notin \Dod(\Fun{h})$,
\begin{equation}
\label{eq:comp-fpr}
\fprate_{\Fun{g} \circ \Fun{f}}
\;=\;
\fprate_1
\bigl[\rho\,(1-\fnrate_2) + (1-\rho)\,\fprate_2\bigr]\,.
\end{equation}
\end{theorem}

\begin{proof}
Fix $x \notin \Dod(\Fun{h}) = \Dod(\Fun{f})$.
A false positive in the composition means
$\APFun{h}(x) \neq \Nothing$ despite $x \notin \Dod(\Fun{h})$.
By \cref{eq:comp-def}, this requires two events in
succession.

\emph{Step 1: $\APFun{f}$ must produce a false positive.}
Since $x \notin \Dod(\Fun{f})$,
$\APFun{f}[\fprate_1][\fnrate_1](x) = \Nothing$ with probability
$1 - \fprate_1$ and produces some $y \in \Set{Y}$ with probability
$\fprate_1$.  If \Nothing is returned, the composition returns
\Nothing (a correct negative), so we require the false positive event
from $\APFun{f}$, which occurs with probability $\fprate_1$.

\emph{Step 2: The false positive value $y$ must survive
$\APFun{g}$.}
Suppose $\APFun{f}[\fprate_1][\fnrate_1](x) = y$ (a false positive).
We condition on the location of $y$ relative to $\Dod(\Fun{g})$.

\begin{itemize}
\item \emph{Case $y \in \Dod(\Fun{g})$} (probability $\rho$).
Since $y$ is a positive element for $\Fun{g}$,
$\APFun{g}[\fprate_2][\fnrate_2](y)$ returns a value in $\Set{Z}$
with probability $1 - \fnrate_2$ and returns \Nothing with probability
$\fnrate_2$.  The composition produces a false positive with probability $1 - \fnrate_2$.

\item \emph{Case $y \notin \Dod(\Fun{g})$} (probability $1-\rho$).
Since $y$ is a negative element for $\Fun{g}$,
$\APFun{g}[\fprate_2][\fnrate_2](y)$ returns a value in $\Set{Z}$
with probability $\fprate_2$ and returns \Nothing with probability
$1 - \fprate_2$.  The composition produces a false positive with
probability $\fprate_2$.
\end{itemize}

\emph{Step 3: Combine by the law of total probability.}
By independence:
\begin{align}
\fprate_{\Fun{g} \circ \Fun{f}}
  &= \fprate_1
     \bigl[\rho\,(1-\fnrate_2) + (1-\rho)\,\fprate_2\bigr]\,.
     \label{eq:fpr-derivation}
\end{align}
\end{proof}

\begin{remark}
The false positive rate of the composition is always bounded above by
$\fprate_1$, since the bracketed factor
$\rho(1-\fnrate_2)+(1-\rho)\fprate_2 \leq 1$.  The second map acts
as an additional filter: a false positive from $\APFun{f}$ must also
survive $\APFun{g}$ to become a false positive in the composition.
\end{remark}

We now record two important special cases.

\begin{corollary}[Disjoint false positives]
\label{cor:disjoint-fp}
If $\rho = 0$, i.e., false positives from
$\APFun{f}[\fprate_1][\fnrate_1]$ always land outside
$\Dod(\Fun{g})$, then
\begin{equation}
\label{eq:fpr-product}
\fprate_{\Fun{g} \circ \Fun{f}} = \fprate_1 \cdot \fprate_2\,.
\end{equation}
\end{corollary}
\begin{proof}
Setting $\rho = 0$ in \cref{eq:comp-fpr}:
$\fprate_{\Fun{g} \circ \Fun{f}}
= \fprate_1 [0 \cdot (1-\fnrate_2) + 1 \cdot \fprate_2]
= \fprate_1 \, \fprate_2$.
\end{proof}

This is the ``independent errors compound'' case: the two false positive
mechanisms act in series and the overall rate is the product.

\begin{corollary}[Complete false-positive overlap]
\label{cor:full-overlap}
If $\rho = 1$, i.e., false positives from
$\APFun{f}[\fprate_1][\fnrate_1]$ always land inside
$\Dod(\Fun{g})$, then
\begin{equation}
\label{eq:fpr-full-overlap}
\fprate_{\Fun{g} \circ \Fun{f}} = \fprate_1 \,(1-\fnrate_2)\,.
\end{equation}
\end{corollary}
\begin{proof}
Setting $\rho = 1$ in \cref{eq:comp-fpr}:
$\fprate_{\Fun{g} \circ \Fun{f}}
= \fprate_1 [1 \cdot (1-\fnrate_2) + 0 \cdot \fprate_2]
= \fprate_1 (1-\fnrate_2)$.
\end{proof}

\begin{corollary}[$k$-fold composition]
\label{cor:k-fold}
Let
$\APFun{f}[\fprate][\fnrate]_1, \ldots,
 \APFun{f}[\fprate][\fnrate]_k$
be $k$ independent approximate maps, each with the same rates
$(\fprate, \fnrate)$, and suppose $\rho = 0$ at every stage.
Then the $k$-fold composition has rates
\begin{equation}
\label{eq:k-fold}
\fprate^{(k)} = \fprate^k
\qquad\text{and}\qquad
\fnrate^{(k)} = 1 - (1-\fnrate)^k\,.
\end{equation}
\end{corollary}
\begin{proof}
By induction on $k$.  The base case $k=1$ is immediate.
For the inductive step, applying \cref{cor:disjoint-fp}
gives
$\fprate^{(k)} = \fprate^{(k-1)} \cdot \fprate = \fprate^k$,
and applying \cref{thm:comp-fnr} gives
$1 - \fnrate^{(k)} = (1-\fnrate^{(k-1)})(1-\fnrate) = (1-\fnrate)^k$.
\end{proof}

\begin{remark}
The $k$-fold false positive rate $\fprate^k \to 0$ exponentially as
$k \to \infty$, while the false negative rate
$1-(1-\fnrate)^k \to 1$.  A long composition chain is therefore
increasingly specific (few false positives) but increasingly
insensitive (many false negatives).  This is the fundamental
precision--recall trade-off of iterated approximate computation.
\end{remark}

\begin{example}
\label{ex:two-stage}
Let $\Fun{f} \colon \Set{X} \pfun \Set{Y}$ and
$\Fun{g} \colon \Set{Y} \pfun \Set{Z}$ with
$\fprate_1 = 0.01$, $\fnrate_1 = 0$,
$\fprate_2 = 0.02$, $\fnrate_2 = 0$,
and suppose false positives from $\APFun{f}$ are uniform on $\Set{Y}$
with $\Card{\Dod(\Fun{g})}/\Card{\Set{Y}} = 0.1$,
so $\rho = 0.1$.

By \cref{thm:comp-fnr},
$\fnrate_{\Fun{g} \circ \Fun{f}} = 0$.

By \cref{thm:comp-fpr},
\begin{equation*}
\fprate_{\Fun{g} \circ \Fun{f}}
  = 0.01 \bigl[0.1 \cdot 1 + 0.9 \cdot 0.02\bigr]
  = 0.01 \cdot 0.118
  = 0.00118\,.
\end{equation*}
The composition has a false positive rate roughly one order of
magnitude smaller than either stage alone.
\end{example}

% -------------------------------------------------------------------
\subsection{Higher-order structure and special cases}
\label{sec:higher-order-comp}

\subsubsection{Composition raises the model order}

The first-order approximate map model assumes that all negative
elements are identically distributed.
Composition breaks this symmetry.

\begin{lemma}[Composition raises the model order]
\label{prop:higher-order}
Let $\APFun{f}[\fprate_1][\fnrate_1]$ and
$\APFun{g}[\fprate_2][\fnrate_2]$ be first-order approximate maps.
Then the composition
$\APFun{g}[\fprate_2][\fnrate_2] \circ
 \APFun{f}[\fprate_1][\fnrate_1]$
is, in general, a second-order approximate map.
\end{lemma}
\begin{proof}
Let $x \notin \Dod(\Fun{f})$ and suppose
$\APFun{f}[\fprate_1][\fnrate_1](x) = y$ (a false positive).
The probability that the composition returns a non-\Nothing value
depends on whether $y \in \Dod(\Fun{g})$ or $y \notin \Dod(\Fun{g})$.

For elements whose false positives land in $\Dod(\Fun{g})$, the effective
false positive rate through the composition is
$\fprate_1(1-\fnrate_2)$.
For elements whose false positives land outside $\Dod(\Fun{g})$, it is
$\fprate_1 \cdot \fprate_2$.
Since $(1-\fnrate_2) \neq \fprate_2$ in general, these two subsets have
distinct false positive probabilities, so the composition is not a
first-order model.
\end{proof}

\begin{remark}
If the false positive distribution of $\APFun{f}$ is uniform over
$\Set{Y}$, then for any given $x \notin \Dod(\Fun{f})$, the value of
$\rho$ is the same.  In this case, the negative set is still
identically distributed \emph{marginally} (after integrating over the
random false positive value), and the composition admits a
first-order description with the aggregate rate given by
\cref{eq:comp-fpr}.  The second-order structure is exposed only when
the false positive distribution of $\APFun{f}$ depends on $x$.
\end{remark}

\subsubsection{Set operations as compositions of indicator functions}

The classical set-algebraic operations can be expressed as compositions involving indicator functions.
Let $\Set{A}, \Set{B} \subseteq \Set{X}$ be sets with indicator
functions $\SetIndicator{\Set{A}}, \SetIndicator{\Set{B}} \colon
\Set{X} \to \BitSet$.

\begin{enumerate}
\item \emph{Complement.}
The indicator of $\SetComplement[\Set{A}]$ is
$\SetIndicator{\SetComplement[\Set{A}]}(x) = 1 -
\SetIndicator{\Set{A}}(x)$.
Under the approximate model, negation swaps the roles of positive and
negative, so for an approximate set $\ASet{A}[\fprate][\fnrate]$, the
approximate complement has rates $\fprate' = \fnrate$ and
$\fnrate' = \fprate$.

\item \emph{Intersection.}
The indicator of $\Set{A} \cap \Set{B}$ is
$\SetIndicator{\Set{A} \cap \Set{B}}(x) =
\SetIndicator{\Set{A}}(x) \cdot \SetIndicator{\Set{B}}(x)$.
Under independence,
$\tprate_{\Set{A} \cap \Set{B}} = \tprate_{\Set{A}} \cdot
\tprate_{\Set{B}}$.

\item \emph{Union.}
By De Morgan's law,
$\Set{A} \cup \Set{B} =
\SetComplement[\SetComplement[\Set{A}] \cap
\SetComplement[\Set{B}]]$,
so the union error rates follow from the complement and intersection
formulas.
\end{enumerate}

These results are developed systematically in the companion paper on
composition of Bernoulli sets.  The key observation is that
\cref{thm:comp-fnr} generalizes the intersection formula: the
true-positive-rate product rule for set intersection is the special
case where both maps are indicator functions and the composition is
pointwise multiplication.

\subsubsection{The Boolean function space}

\begin{example}[Boolean maps and approximate NOT]
\label{ex:bool}
Consider the function space
$\Bool \to \Bool = \{\True,\False\} \to \{\True,\False\}$.
There are exactly $2^2 = 4$ total functions:
\begin{center}
\begin{tabular}{@{} l c c @{}}
\toprule
Function & $\True \mapsto$ & $\False \mapsto$ \\
\midrule
$\Fun{id}$          & $\True$  & $\False$ \\
$\Fun{not}$         & $\False$ & $\True$  \\
$\Fun{const\text{-}T}$ & $\True$  & $\True$  \\
$\Fun{const\text{-}F}$ & $\False$ & $\False$ \\
\bottomrule
\end{tabular}
\end{center}

Each of these is a \emph{total} function, so there are no domain-level false
positives: $\fprate = 0$.
Consider the approximate negation
$\APFun{not}[\fprate][\fnrate]$ with $\fprate = 0$ and false negative rate $\fnrate$.
We compute $\APFun{not}[\fprate][\fnrate] \circ
\APFun{not}[\fprate][\fnrate]$, which approximates
$\Fun{not} \circ \Fun{not} = \Fun{id}$.

By \cref{thm:comp-fnr} with $\fnrate_1 = \fnrate_2 = \fnrate$:
\begin{equation}
\fnrate_{\Fun{id}} = 1 - (1-\fnrate)^2 = 2\fnrate - \fnrate^2\,.
\end{equation}
For small $\fnrate$, the composed false negative rate is approximately
$2\fnrate$: each stage independently
contributes an $\fnrate$ probability of failure.
For instance, with $\fnrate = 0.05$, the two-fold composition has
$\fnrate_{\Fun{id}} = 1 - 0.95^2 = 0.0975$.
\end{example}

\end{document}
